\section*{अध्याय \ref{ch:g6:rocks-and-minerals} --- चट्टानें और खनिज}

\begin{solution}{exo:g6:rocks-and-minerals:1}
ग्रेनाइट: आग्नेय। बलुआ पत्थर: अवसादी। संगमरमर: कायांतरित। बेसाल्ट:
आग्नेय। चूना पत्थर: अवसादी। स्लेट: कायांतरित।
\end{solution}

\begin{solution}{exo:g6:rocks-and-minerals:2}
क्वार्ट्ज़ एक \omterm{def:g6:rocks-and-minerals:mineral}{खनिज} है; ग्रेनाइट एक \omterm{def:g3:stones-pebbles-sand:rock}{चट्टान} है। \omterm{def:g6:rocks-and-minerals:mineral}{खनिज} एक अकेला प्राकृतिक पदार्थ है
जिसके परमाणु एक नियमित ढंग से सजे होते हैं; \omterm{def:g3:stones-pebbles-sand:rock}{चट्टान} आपस में चिपके एक या कई \omterm{def:g6:rocks-and-minerals:mineral}{खनिजों} से
बनी होती है (ग्रेनाइट: फ़ेल्डस्पार, क्वार्ट्ज़ और अभ्रक)।
\end{solution}

\begin{solution}{exo:g6:rocks-and-minerals:3}
टैल्क (1), कैल्साइट (3), क्वार्ट्ज़ (7), हीरा (10)। हीरा, जो सबसे कठोर है,
बाक़ी सबको खरोंचता है।
\end{solution}

\begin{solution}{exo:g6:rocks-and-minerals:4}
``पिघली \omterm{def:g3:stones-pebbles-sand:rock}{चट्टान} ठंडा होने से \omterm{def:g6:rocks-and-minerals:rock-families}{आग्नेय चट्टान} बनती है। चूना पत्थर ताप और दाब से
संगमरमर बनता है। ग्रेनाइट \omterm{def:g5:water-shapes-the-land:weathering}{अपक्षय} से \omterm{def:g3:stones-pebbles-sand:sand}{रेत} के कणों में
बदलता है।''
\end{solution}

\begin{solution}{exo:g6:rocks-and-minerals:5}
(क) बेसाल्ट। (ख) संगमरमर। (ग) बलुआ पत्थर।
\end{solution}

\begin{solution}{exo:g6:rocks-and-minerals:6}
उसकी \omterm{def:g6:rocks-and-minerals:hardness}{कठोरता} नाख़ून की \omterm{def:g6:rocks-and-minerals:hardness}{कठोरता} (लगभग 2.5) से ज़्यादा और ताँबे के सिक्के की \omterm{def:g6:rocks-and-minerals:hardness}{कठोरता}
(लगभग 3.5) से कम है। वह कैल्साइट (3) हो सकता है।
\end{solution}

\begin{solution}{exo:g6:rocks-and-minerals:7}
क्वार्ट्ज़: $21 \div 60 = 0.35$, यानी 35~\%। फ़ेल्डस्पार: $33 \div 60 = 0.55$,
यानी 55~\%। अभ्रक: $6 \div 60 = 0.10$, यानी 10~\%। जाँच:
$35 + 55 + 10 = 100$।
\end{solution}

\begin{solution}{exo:g6:rocks-and-minerals:8}
ग्रेनाइट ज़मीन के बहुत नीचे बहुत धीरे ठंडा होता है: उसके कणों को बढ़ने का लंबा समय
मिलता है, और वे बड़े हो जाते हैं। बेसाल्ट \omterm{def:g4:volcanoes-earthquakes-first-look:eruption}{लावा} से बनता है, जो सतह पर जल्दी ठंडा
होता है: उसके कणों को बढ़ने का समय नहीं मिलता और वे देखने के लिए बहुत छोटे रहते हैं।
\end{solution}

\begin{solution}{exo:g6:rocks-and-minerals:9}
(क) चूना पत्थर अवसादी है; संगमरमर कायांतरित है। (ख) ज़मीन के बहुत नीचे
ताप और दाब ने उसे बिना पिघलाए बदल दिया: उसका कैल्साइट बढ़कर आपस में फ़िट बैठने वाले
चमकीले क्रिस्टल बन गया। (ग) दोनों एक ही \omterm{def:g6:rocks-and-minerals:mineral}{खनिज}, कैल्साइट, से बने हैं, और
कैल्साइट ही अम्ल से बुदबुदाता है।
\end{solution}

\begin{solution}{exo:g6:rocks-and-minerals:10}
(क) हाँ: क्वार्ट्ज़ (7) काँच (लगभग 5.5) से कठोर है। (ख) नहीं: फ़्लोराइट (4)
काँच से नरम है। (ग) हाँ: स्टील की कील (लगभग 6.5) फ़ेल्डस्पार (6) से
कठोर है।
\end{solution}

\begin{solution}{exo:g6:rocks-and-minerals:11}
उदाहरण के लिए: (1) \omterm{def:g1:landscapes-around-us:landform}{पर्वत} का ग्रेनाइट (आग्नेय) \omterm{def:g5:water-shapes-the-land:weathering}{अपक्षय} से टूटता है, और उसके
क्वार्ट्ज़ कण आज़ाद हो जाते हैं; (2) एक \omterm{def:g1:landscapes-around-us:water}{नदी} उन्हें बहा ले जाती है:
\omterm{def:g5:water-shapes-the-land:erosion}{अपरदन}; (3) कण \omterm{def:g2:water-on-earth:ocean}{समुद्र} की तली पर परतों में बिछ जाते हैं:
\omterm{def:g5:water-shapes-the-land:deposition}{निक्षेपण}; (4) दूसरी परतों के नीचे दबकर वे दबाए और जोड़े जाकर बलुआ पत्थर बन
जाते हैं: कठोर होना (\omterm{def:g3:stones-pebbles-sand:rock}{चट्टान} अब अवसादी है);
(5) ज़मीन के बहुत नीचे ले जाए जाने पर बलुआ पत्थर ताप और दाब से एक \omterm{def:g6:rocks-and-minerals:rock-families}{कायांतरित
चट्टान} में बदल जाता है।
\end{solution}

\begin{solution}{exo:g6:rocks-and-minerals:12}
(क) स्टील (लगभग 6.5) से नहीं खरोंचा जाता और काँच को खरोंच सकता है: 6.5 से
ज़्यादा कठोर। पुखराज (8) से खरोंच जाता है: 8 से नरम। उसकी \omterm{def:g6:rocks-and-minerals:hardness}{कठोरता} 7 है: वह
क्वार्ट्ज़ है। (ख) क्वार्ट्ज़ उससे रगड़ने वाली लगभग हर चीज़ (दूसरे कण, \omterm{def:g3:stones-pebbles-sand:pebble}{कंकड़}) से कठोर
है: वह मुश्किल से घिसता है। कैल्साइट (3) और अभ्रक
(2 से 3) बहुत नरम हैं: रगड़ उन्हें पीसकर नन्हे टुकड़ों में बदल देती है जो
बह जाते हैं।
\end{solution}

\begin{solution}{pb:g6:rocks-and-minerals:1}
\textbf{1.} ग्रेनाइट, एक \omterm{def:g6:rocks-and-minerals:rock-families}{आग्नेय चट्टान}; उसके \omterm{def:g6:rocks-and-minerals:mineral}{खनिज} फ़ेल्डस्पार (गुलाबी या
सफ़ेद), क्वार्ट्ज़ (धूसर, काँच जैसा) और अभ्रक (काला, चमकीला) हैं।

\textbf{2.} चूना पत्थर, एक \omterm{def:g6:rocks-and-minerals:rock-families}{अवसादी चट्टान}। खोल उन जानवरों के थे
जो \omterm{def:g2:water-on-earth:ocean}{समुद्र} में रहते थे; उनके अवशेष बिछे और कठोर होकर
\omterm{def:g3:stones-pebbles-sand:rock}{चट्टान} बन गए।

\textbf{3.} स्लेट, एक \omterm{def:g6:rocks-and-minerals:rock-families}{कायांतरित चट्टान}, जो शेल से बनी।

\textbf{4.} संगमरमर, एक \omterm{def:g6:rocks-and-minerals:rock-families}{कायांतरित चट्टान}, जो चूना पत्थर से बनी।

\textbf{5.} क्वार्ट्ज़: स्टील (लगभग 6.5) से नहीं खरोंचा जाता, वह काँच
(लगभग 5.5) को खरोंचता है; उसकी \omterm{def:g6:rocks-and-minerals:hardness}{कठोरता} 7 है।

\textbf{6.} स्टील (लगभग 6.5) से खरोंच जाता है पर काँच (लगभग 5.5) से नहीं:
5.5 और 6.5 के बीच। वह फ़ेल्डस्पार (ऑर्थोक्लेज़, 6) है।

\textbf{7.} अभ्रक (\omterm{def:g6:rocks-and-minerals:hardness}{कठोरता} 2 से 3)।

\textbf{8.} हाँ: फ़र्श चूना पत्थर का है, जो कैल्साइट (3) से बना है, जो ताँबे के
सिक्के (लगभग 3.5) से नरम है।

\textbf{9.} \omterm{def:g5:water-shapes-the-land:weathering}{अपक्षय}: \omterm{def:g4:clouds-rain-snow:precipitation}{बारिश} पत्थर की सतह को धीरे-धीरे घिस
देती है।

\textbf{10.} \omterm{def:g1:landscapes-around-us:landform}{पर्वत} का \omterm{def:g5:water-shapes-the-land:weathering}{अपक्षय} होता है; उसके क्वार्ट्ज़ कण \omterm{def:g1:landscapes-around-us:water}{नदियाँ} ले जाती हैं और
किसी \omterm{def:g1:landscapes-around-us:water}{झील} या \omterm{def:g2:water-on-earth:ocean}{समुद्र} की तली पर परतों में बिछा देती हैं; दबकर
और कठोर होकर वे बलुआ पत्थर, एक \omterm{def:g6:rocks-and-minerals:rock-families}{अवसादी चट्टान}, बन सकते हैं।

\textbf{11.} तीनों परिवार मौजूद हैं: आग्नेय (फ़व्वारे का
ग्रेनाइट), अवसादी (फ़र्श का चूना पत्थर), कायांतरित (छत की
स्लेट और मूर्ति का संगमरमर)।

\textbf{12.} सबसे कठोर \omterm{def:g6:rocks-and-minerals:mineral}{खनिज} ग्रेनाइट का क्वार्ट्ज़ है: मोह्स
\omterm{def:g6:rocks-and-minerals:hardness}{कठोरता} \textbf{7}।
\end{solution}
